% Extracción quirúrgica del propietario V2 05e: líneas 68--110.
% El resto permanece en este archivo como procedencia, pero no se publica.
\iffalse
\chapter[Internal Determinantal Macrostructure]{Internal Determinantal Macrostructure Generated by the Discrete Continuum}
\label{chap:detint-macroestructura}

\begin{thesisbox}
The \({\rm det}\) branch starts exclusively from surviving histories of
\(\mathcal C_{\rm cont}^{\rm disc}\). The nerve of extensions, the Alexander--Whitney diagonal,
the local complex, the unit, the two publications, the filler, the terminal
reduction, and the \(3\)-adic tower are constructed before any subsequent
evaluation. The balance and acyclicity conditions are incorporated into the
typed domain and are never tacitly assumed.
\end{thesisbox}

\section{Surviving determinant domain}

Let
\begin{equation}\label{eq:detint-an}
 A_n:=\mathbf Z/3^n\mathbf Z,
 \qquad
 \rho_{n+1,n}:A_{n+1}\longrightarrow A_n.
\end{equation}
For \(x\in\mathcal C_{\rm cont}^{\rm disc}\), let
\(\mathsf T_m(x)\) be the finite category of TPK prefixes of depth \(m\)
compatible with \(x\). Each object preserves sheet, orientation, residue,
quotient, integer carry, memory, boundary, route, multisection, survival,
and terminal; its morphisms are the admissible extensions that preserve those
coordinates.

From each object \(y\in\mathsf T_m(x)\), the internal integers are extracted
\begin{equation}\label{eq:detint-c-v-lifts}
 \widetilde c_y,\widetilde v_y\in\mathbf Z,
 \qquad
 c_{y,n}:=\widetilde c_y\bmod3^n,
 \qquad
 v_{y,n}:=\widetilde v_y\bmod3^n.
\end{equation}
The first coordinate is the unoriented carry and the second, the oriented
boundary coefficient. Reversal satisfies
\begin{equation}\label{eq:detint-or-cv}
 c_{\bar y,n}=c_{y,n},
 \qquad
 v_{\bar y,n}=-v_{y,n}.
\end{equation}

We write
\begin{equation}\label{eq:detint-z3-k3}
 \mathbf Z_3:=\varprojlim_n A_n,
 \qquad
 K_3:=\operatorname{Frac}(\mathbf Z_3).
\end{equation}
The determinantial restriction is defined by
\begin{equation}\label{eq:detint-dominio}
\begin{split}
 \mathcal C_{\rm cont}^{\rm det}:=
 \bigl\{x\in\mathcal C_{\rm cont}^{\rm disc}:{}&
 x\text{ survives at every depth};\\
 &\widehat P_x^{\rm red}\text{ is finite and based};\\
 &\chi(\widehat P_x^{\rm red})=0;\\
 &H_*(\widehat P_x^{\rm red}\otimes_{\mathbf Z_3}K_3)=0;\\
 &\text{the canonical reduction produces a square block}\\
 &\text{stable under refinement}\bigr\}.
\end{split}
\end{equation}
The objects \(\widehat P_x^{\rm red}\) and the square block will be constructed
in \cref{sec:detint-terminal}. Therefore
\eqref{eq:detint-dominio} is a verifiable membership condition and not
a conclusion inserted into the definition of the maps.

\fi
\section{Nerve of routes, Alexander--Whitney diagonal and counit}

The normalized complex is defined
\begin{equation}\label{eq:detint-gmn}
 G_{m,n}(x):=
 N_*^{\rm norm}\!\left(N\mathsf T_m(x);A_n\right).
\end{equation}
For a nondegenerate simplex
\(\sigma=[y_0\to y_1\to\cdots\to y_q]\), the diagonal is
\begin{equation}\label{eq:detint-aw}
 \Delta_{\rm AW}(\sigma)
 =\sum_{j=0}^q
 [y_0\to\cdots\to y_j]\otimes
 [y_j\to\cdots\to y_q].
\end{equation}
The counit is given by
\begin{equation}\label{eq:detint-counidad-g}
 \epsilon_G([y])=1,
 \qquad
 \epsilon_G(\sigma)=0\quad(q>0).
\end{equation}
In particular,
\begin{equation}\label{eq:detint-vertice-grouplike}
 \Delta_{\rm AW}[y]=[y]\otimes[y],
 \qquad
 \epsilon_G([y])=1.
\end{equation}

If \(m'\ge m\) and \(n'\ge n\), truncation and coefficient reduction
induce
\begin{equation}\label{eq:detint-rmn}
 R_{m',m}^{n',n}:G_{m',n'}(x)\longrightarrow G_{m,n}(x),
\end{equation}
and they are verified
\begin{equation}\label{eq:detint-aw-natural}
 R\partial=\partial R,
 \qquad
 (R\otimes R)\Delta_{\rm AW}=\Delta_{\rm AW}R,
 \qquad
 \epsilon_GR=\rho_{n',n}\epsilon_G.
\end{equation}
The path is not reduced to its endpoints: the entire simplex remains in
\(G_{m,n}(x)\).

\iffalse
\section{Internal local complex}

For \(c\in A_n\), define
\begin{equation}\label{eq:detint-p0-grados}
 P^0_{n,1}(c)=A_nf,
 \qquad
 P^0_{n,0}(c)=A_nr\oplus A_ne\oplus A_nb,
 \qquad
 P^0_{n,-1}(c)=A_ng,
\end{equation}
with differential
\begin{equation}\label{eq:detint-p0-diferencial}
 \partial f=3ce+b,
 \qquad
 \partial r=0,
 \qquad
 \partial e=g,
 \qquad
 \partial b=-3cg.
\end{equation}
The chain identity is literal:
\begin{equation}\label{eq:detint-d2}
 \partial^2f=3c\,\partial e+\partial b=3cg-3cg=0,
 \qquad
 \partial^2e=\partial^2b=\partial^2r=0.
\end{equation}

The increase
\begin{equation}\label{eq:detint-epsilon-p}
 \epsilon_P:P_n^0(c)\longrightarrow A_n[0]
\end{equation}
is fixed by
\[
 \epsilon_P(r)=1,
 \qquad
 \epsilon_P(e)=\epsilon_P(b)=\epsilon_P(f)=\epsilon_P(g)=0.
\]
The complex possesses the differential coalgebra.
\begin{equation}\label{eq:detint-coalgebra-p}
 \Delta_P(r)=r\otimes r,
 \qquad
 \Delta_P(q)=q\otimes r+r\otimes q
 \quad(q\in\{e,b,f,g\}).
\end{equation}
Thus, \(r\) is group-like and the other four generators are primitive
relative to \(r\). Equations
\eqref{eq:detint-p0-diferencial} prove that \(\Delta_P\) commutes with the
tensor differential.

\section{Local extension system and memory ledger}

For an extension \(\alpha:y\to y'\), we define
\begin{equation}\label{eq:detint-psi}
 \Psi_\alpha:P_n^0(c_y)\longrightarrow P_n^0(c_{y'})
\end{equation}
through means of
\begin{equation}\label{eq:detint-psi-generadores}
\begin{gathered}
 \Psi_\alpha(r_y)=r_{y'},\quad
 \Psi_\alpha(e_y)=e_{y'},\quad
 \Psi_\alpha(g_y)=g_{y'},\quad
 \Psi_\alpha(f_y)=f_{y'},\\
 \Psi_\alpha(b_y)=b_{y'}+3(c_{y'}-c_y)e_{y'}.
\end{gathered}
\end{equation}
It is a map of complexes because
\begin{equation}\label{eq:detint-psi-b-chain}
 \partial\Psi_\alpha(b_y)
 =-3c_{y'}g+3(c_{y'}-c_y)g
 =-3c_yg
 =\Psi_\alpha(\partial b_y),
\end{equation}
and
\begin{equation}\label{eq:detint-psi-f-chain}
 \Psi_\alpha(\partial f_y)
 =3c_ye+b+3(c_{y'}-c_y)e
 =3c_{y'}e+b
 =\partial\Psi_\alpha(f_y).
\end{equation}
Moreover,
\begin{equation}\label{eq:detint-psi-functorial}
 \Psi_\beta\Psi_\alpha=\Psi_{\beta\alpha}.
\end{equation}

The oriented coordinate change does not erase. It is recorded by
\begin{equation}\label{eq:detint-k-alpha}
 K_\alpha:=(v_{y'}-v_y)f_{y'}.
\end{equation}
For composable extensions,
\begin{equation}\label{eq:detint-k-coherencia}
 K_{\beta\alpha}=K_\beta+\Psi_\beta K_\alpha.
\end{equation}
This equality is the exact ledger of the boundary increment and makes visible
the memory that a purely nominal composition would lose.

The complete local diagram is the Grothendieck construction.
\begin{equation}\label{eq:detint-grothendieck-local}
 \int_{y\in\mathsf T_m(x)}P_n^0(c_y).
\end{equation}
In it, an arrow over \(\alpha:y\to y'\) carries \(\Psi_\alpha\)
as its linear component and \(K_\alpha\) as its coherence. When applied to an
element supported at a vertex, the notation is used
\begin{equation}\label{eq:detint-psi-soporte-vertice}
 \Psi_\alpha(p,[y]):=(\Psi_\alpha p,[y']).
\end{equation}
The edge \([y\to y']\) remains in the path coordinate of the nerve. This
convention will be used in the naturality formulas for
\(z_\pm\) and \(h_*\); no global map erasing the remaining
simplices of \(G_{m,n}(x)\) is asserted.

\section{Pullback, unit, root, publications and filler}

The pullback of complexes is
\begin{equation}\label{eq:detint-pgen}
 P^{\rm gen}_{y;m,n}
 :=P_n^0(c_y)\times_{A_n[0]}G_{m,n}(x)
 =\{(p,\gamma):\epsilon_P(p)=\epsilon_G(\gamma)\}.
\end{equation}
Its differential is
\begin{equation}\label{eq:detint-pgen-d}
 \partial(p,\gamma)=(\partial p,\partial\gamma).
\end{equation}
The common unit associated with the vertex \(y\) is
\begin{equation}\label{eq:detint-unidad}
 \mathbf1_y:=(r,[y]).
\end{equation}
Its two projections are group-like by
\eqref{eq:detint-vertice-grouplike} and
\eqref{eq:detint-coalgebra-p}; that is the precise meaning of a compatible
group-like unit of the pullback.

The root projection is the map of complexes
\begin{equation}\label{eq:detint-proyeccion-root}
 \varpi_{\rm root}:P^{\rm gen}_{y;m,n}\longrightarrow
 \operatorname{Root}_{y;m,n},
 \qquad
 \varpi_{\rm root}(p,\gamma):=(\epsilon_P(p)r,\gamma).
\end{equation}
The two internal publications are defined.
\begin{equation}\label{eq:detint-zpm}
 z_-(y):=(r,[y]),
 \qquad
 z_+(y):=(r+3c_yv_ye+v_yb,[y]).
\end{equation}
Both are cycles:
\begin{equation}\label{eq:detint-zpm-ciclos}
 \partial z_-(y)=0,
 \qquad
 \partial z_+(y)=3c_yv_yg-3c_yv_yg=0.
\end{equation}
Their roots coincide:
\begin{equation}\label{eq:detint-raiz-comun}
 \varpi_{\rm root}z_-(y)
 =(r,[y])
 =\varpi_{\rm root}z_+(y).
\end{equation}

The filler is constructed directly:
\begin{equation}\label{eq:detint-hstar}
 h_*(y):=(-v_yf,0)\in P^{\rm gen}_{y;m,n,1}.
\end{equation}
Then
\begin{equation}\label{eq:detint-filler-identidad}
 \partial h_*(y)
 =(-3c_yv_ye-v_yb,0)
 =z_-(y)-z_+(y).
\end{equation}
In particular,
\begin{equation}\label{eq:detint-clases-iguales}
 [z_-(y)]=[z_+(y)]
 \quad\text{in}\quad H_0(P^{\rm gen}_{y;m,n}),
\end{equation}
and the equality preserves its witness chain.

Under an extension \(\alpha:y\to y'\), we have
\begin{equation}\label{eq:detint-z-natural}
 z_-(y')=\Psi_\alpha z_-(y),
 \qquad
 z_+(y')-\Psi_\alpha z_+(y)=\partial K_\alpha,
\end{equation}
and
\begin{equation}\label{eq:detint-h-natural}
 h_*(y')-\Psi_\alpha h_*(y)=-K_\alpha.
\end{equation}
Therefore, naturality does not identify distinct values of \(v\): it records
their difference through the coherent homotopy \(K_\alpha\).

\section{Internal Orientation}

On \(P_n^0(c)\) is defined
\begin{equation}\label{eq:detint-omega-p}
 \Omega(r)=r,
 \qquad
 \Omega(q)=-q\quad(q\in\{e,b,f,g\}).
\end{equation}
It is an automorphism of complexes and of the relative coalgebra. In the nerve,
the oriented inversion is
\begin{equation}\label{eq:detint-omega-g}
 \mathfrak o_G[y_0\to\cdots\to y_q]
 :=(-1)^{q(q+1)/2}
 [\bar y_q\to\cdots\to\bar y_0].
\end{equation}
Together with \eqref{eq:detint-or-cv}, these actions give
\begin{equation}\label{eq:detint-or-z-h}
 \Omega z_-(y)=z_-(\bar y),
 \qquad
 \Omega z_+(c_y,v_y)=z_+(c_{\bar y},v_{\bar y}),
 \qquad
 \Omega h_*(y)=h_*(\bar y).
\end{equation}
Orientation acts on path, complex, publications, filler and determinant
line; it is not a label without an action.

\section{Terminal reduction and survival condition}
\label{sec:detint-terminal}

Only the common root is eliminated:
\begin{equation}\label{eq:detint-p-reducido}
 \overline P_{y;m,n}
 :=P^{\rm gen}_{y;m,n}/A_n\mathbf1_y,
 \qquad
 \widehat P_{y;m}^{\rm red}:=\varprojlim_n\overline P_{y;m,n}.
\end{equation}
The base is first ordered by the TPK order of routes and then by
\begin{equation}\label{eq:detint-orden-base}
 f\prec e\prec b\prec g.
\end{equation}
The following internal algorithm is executed.
\begin{enumerate}
 \item Find the coefficients that are units of \(\mathbf Z_3\).
 \item Choose the first according to \eqref{eq:detint-orden-base} and the order of
 paths.
 \item Cancel the corresponding pair and record the elementary
 operations, their sign, and their contribution to the orientation.
 \item Repeat until no cancellable unit pivot remains.
\end{enumerate}
The survival condition of \eqref{eq:detint-dominio} requires that,
cofinally in \(m\), the result be a two-term block
\begin{equation}\label{eq:detint-sx}
 \mathbf Z_3^{r_x}\xrightarrow{S_x}\mathbf Z_3^{r_x},
 \qquad
 \det S_x\ne0\text{ in }K_3,
\end{equation}
and that a refinement only add unit contractible blocks and recorded
changes of basis. If these conditions fail, the history does not belong to
\(\mathcal C_{\rm cont}^{\rm det}\); a nonexistent terminal is not forced.

The identity route provides an elementary witness. After removing the root,
the local complex is
\begin{equation}\label{eq:detint-testigo-contractil}
 A_nf\xrightarrow{\binom{3c}{1}}
 A_ne\oplus A_nb\xrightarrow{(1,-3c)}A_ng.
\end{equation}
Is contractible via
\begin{equation}\label{eq:detint-contraccion-local}
 s(g)=e,
 \qquad
 s(ae+\beta b)=\beta f,
 \qquad
 s(f)=0.
\end{equation}
Thus, the terminal conditions admit at least the internal identity witness
and do not describe a formally empty subdomain.

\section{Smith, \texorpdfstring{\(3\)}{3}-adic order and initial unit}

In oriented bases, a valuation diagonalization has the form
\begin{equation}\label{eq:detint-smith}
 U_xS_xV_x
 =\operatorname{diag}(3^{a_1}u_1,\ldots,3^{a_r}u_r),
 \qquad
 a_i\in\mathbf N,quad u_i\in\mathbf Z_3^\times.
\end{equation}
We define
\begin{equation}\label{eq:detint-orden-d}
 d_x:=v_3(\det S_x)=\sum_{i=1}^r a_i
\end{equation}
and the initial unit
\begin{equation}\label{eq:detint-leading-unit}
 \lambda_x:=3^{-d_x}\det S_x\in\mathbf Z_3^\times.
\end{equation}
The orientation coordinate trivializes the determinant line. If a
transition produces
\begin{equation}\label{eq:detint-estabilidad-bloques}
 S_x\oplus I_q=U S_{x'}V,
\end{equation}
the orientation is transported by the inverse factor of \(\det U\det V\).
The oriented unit \(\lambda_x^{\rm or}\) is then natural and does not change
when contractible blocks are added.

At finite precision,
\begin{equation}\label{eq:detint-smith-finito}
 S_{x,n}=S_x\bmod3^n,
 \qquad
 a_i^{(n)}=\min(a_i,n).
\end{equation}
The terminal is, therefore, the compatible family of all precisions
and not an isolated integer.

\section{Bockstein Tower and Higher Order Carry}

Let
\[
 C_x^1\xrightarrow{S_x}C_x^0
\]
the terminal block. If a class \([\bar y]\) admits a lift
\(y\in C_x^1\) with \(S_xy\in3^qC_x^0\), we define
\begin{equation}\label{eq:detint-beta-q}
 \beta_q[\bar y]
 :=\left[3^{-q}S_xy\bmod3\right].
\end{equation}
Changing the lift alters the representative by a boundary of the corresponding
page, so \(\beta_q\) is well defined. The higher
carry is
\begin{equation}\label{eq:detint-carry-superior}
 \operatorname{Carr}_q(y):=3^{-q}S_xy.
\end{equation}
It does not replace the TPK carry of \eqref{eq:detint-c-v-lifts}; it records its
prolongation in the coefficient tower.

For a block \(3^{a_i}u_i\), the class persists until page \(a_i\),
and on that page the nonzero differential is multiplication by
\(\bar u_i\in\mathbf F_3^\times\). Consequently,
\begin{equation}\label{eq:detint-bock-longitudes}
 \{\text{survival lengths}\}=\{a_1,\ldots,a_r\},
 \qquad
 d_x=\sum_i a_i.
\end{equation}
The oriented product of the page trivializations satisfies
\begin{equation}\label{eq:detint-bock-leading}
 \Theta_{\rm Bock}(x)=\overline{\lambda_x^{\rm or}}
 \in\mathbf F_3^\times.
\end{equation}
The proof is carried out in each diagonal block and transported by the recorded
unimodular changes. Smith, Bockstein and the initial unit are thus three
readings of the same terminal.

\section{The thirteen coordinates of the determinantal constraint}

For each \((m,n)\), it is defined
\begin{equation}\label{eq:detint-d-trece}
 D_{{\rm det},m,n}(x):=
 (\mathcal F,\operatorname{Ext}_\Gamma,\operatorname{Rel},
 \mathcal N,\operatorname{or},\operatorname{Carr},\operatorname{Mem},
 \partial,\gamma,\operatorname{Msect},\operatorname{Surv},
 \operatorname{Term},\mathcal L)_{{\rm det},m,n}(x),
\end{equation}
with the following content.
\begin{enumerate}
 \item The fiber is
 \[
  \mathcal F=(A_n,\mathsf T_m(x),
  \{P_n^0(c_y)\}_y,\{P^{\rm gen}_{y;m,n}\}_y).
 \]
 \item \(\operatorname{Ext}_\Gamma\) contains all morphisms
 \(\alpha\), the maps \(\Psi_\alpha\), the homotopies \(K_\alpha\) and
 the maps \(R_{m',m}^{n',n}\).
 \item \(\operatorname{Rel}\) contains \(\partial^2=0\), the
 pullback equation, AW, counit, \(\Psi_\beta\Psi_\alpha=\Psi_{\beta\alpha}\) and
 \(K_{\beta\alpha}=K_\beta+\Psi_\beta K_\alpha\).
 \item
 \(\mathcal N=(m,n,\widetilde c_y,\widetilde v_y,c_{y,n},v_{y,n},
 r_x,a_1,\ldots,a_{r_x},d_x)\).
 \item \(\operatorname{or}\) contains \(y\mapsto\bar y\),
 \(\mathfrak o_G\), \(\Omega\) and the orientation of the determinant
 line.
 \item \(\operatorname{Carr}\) separately preserves the integer carry,
 its reduction \(c_{y,n}\) and all higher carries
 \(\operatorname{Carr}_q\).
 \item \(\operatorname{Mem}\) contiene
 \[
  y_0\to\cdots\to y_m,
  \quad\{(\widetilde c_{y_j},\widetilde v_{y_j},\mathbf1_{y_j})\}_{j=0}^m,
  \quad\{K_{\alpha_j}\}.
 \]
 \item
 \(\partial=(\partial_G,\partial_P,\partial_{P^{\rm gen}},
 \partial_{\overline P},h_*)\).
 \item \(\gamma=[y_0\to\cdots\to y_q]\) preserves the complete oriented
 simplex.
 \item \(\operatorname{Msect}\) contains all compatible
 extensions, all \(3\)-adic lifts, and all unimodular
 presentations of the same oriented terminal; it does not choose one a posteriori.
 \item \(\operatorname{Surv}\) records arbitrary depth,
 coefficient compatibility, balance, acyclicity over \(K_3\),
 stability under unit blocks and persistence of the root.
 \item
 \[
  \operatorname{Term}=(\widehat P_x^{\rm red},S_x,
  \{a_i,u_i\},d_x,\lambda_x^{\rm or},
  \{\beta_q\},\Theta_{\rm Bock}).
 \]
 \item The internal reader is
 \[
  \mathcal L_{\rm det}^{\rm int}
  =(\mathbf1_y,
  \varpi_{\rm root}z_-=\varpi_{\rm root}z_+,
  [z_-]=[z_+],d_x,\lambda_x^{\rm or},\Theta_{\rm Bock}).
 \]
\end{enumerate}

The transitions satisfy, coordinate by coordinate,
\begin{equation}\label{eq:detint-naturalidad-trece}
 u^{\rm det}_{(m',n'),(m,n)}D_{{\rm det},m',n'}(x')
 =D_{{\rm det},m,n}(\tau_{m',m}x').
\end{equation}

\section{Assembler, publisher, and non-ablative chain}

We define
\begin{equation}\label{eq:detint-graph-d}
 \mathcal G_{\rm det}:=\operatorname{Graph}(D_{\rm det}),
 \qquad
 \operatorname{Res}_{\rm det}(x):=(x,D_{\rm det}x).
\end{equation}
Let \(\chi_{\rm det}(x)\) be the complete multisection of germs
\[
 (y,\widetilde c_y,\widetilde v_y,\operatorname{or}_y,\gamma_y)
\]
compatible at all depths. The dependent object is
\begin{equation}\label{eq:detint-x-dependiente}
 X_{\chi,{\rm det}}
 :=\{(\chi_{\rm det}(x),x,D_{\rm det}x):
 x\in\mathcal C_{\rm cont}^{\rm det}\},
\end{equation}
and
\begin{equation}\label{eq:detint-prx}
 \operatorname{pr}_{X,{\rm det}}(x,D_{\rm det}x)
 :=(\chi_{\rm det}(x),x,D_{\rm det}x).
\end{equation}

The raw assembler
\begin{equation}\label{eq:detint-ensamblador-crudo}
 a_{\rm det}:X_{\chi,{\rm det}}\longrightarrow
 \mathscr M_{\rm det}^{\rm amb}
\end{equation}
yields
\begin{equation}\label{eq:detint-ensamblador-salida}
\begin{split}
 a_{\rm det}(z)=\bigl(&
 \{G_{m,n},\Delta_{\rm AW},\epsilon_G\}_{m,n};\\
 &\{P_n^0(c),\Psi_\alpha,K_\alpha,P^{\rm gen},
 \mathbf1,z_-,z_+,h_*\}_{m,n};\\
 &\{\widehat P^{\rm red},S,d,\lambda^{\rm or},
 \beta_q,\Theta_{\rm Bock}\}\bigr).
\end{split}
\end{equation}
The macroobject preserves its entry:
\begin{equation}\label{eq:detint-graph-a}
 \mathfrak M_{\rm det}:=\operatorname{Graph}(a_{\rm det}),
 \qquad
 \mathcal A_{\rm det}(z):=(z,a_{\rm det}z).
\end{equation}

The raw publisher
\begin{equation}\label{eq:detint-publicador-crudo}
 p_{\rm det}:\mathfrak M_{\rm det}\longrightarrow\mathscr Y_{\rm det}
\end{equation}
publishes the certificate jointly
\begin{equation}\label{eq:detint-certificado-publicado}
\begin{gathered}
 \partial^2=0,quad \text{AW and counit},quad
 \mathbf1_y\text{ compatible group-like},\\
 \partial z_\pm=0,quad
 \varpi_{\rm root}z_-=\varpi_{\rm root}z_+,quad
 \partial h_*=z_--z_+,\\
 \Psi_\beta\Psi_\alpha=\Psi_{\beta\alpha},quad
 K_{\beta\alpha}=K_\beta+\Psi_\beta K_\alpha,\\
 d=v_3(\det S)=\sum_i a_i,quad
 \Theta_{\rm Bock}=\overline{\lambda^{\rm or}},quad
 \text{full naturality under refinement}.
\end{gathered}
\end{equation}
Finally,
\begin{equation}\label{eq:detint-graph-p}
 \mathcal C_{\rm det}^{\rm HMT}:=\operatorname{Graph}(p_{\rm det}),
 \qquad
 \mathcal P_{\rm det}(M):=(M,p_{\rm det}M).
\end{equation}
The complete chain is
\begin{equation}\label{eq:detint-cadena-completa}
 \mathcal C_{\rm cont}^{\rm disc}\supset
 \mathcal C_{\rm cont}^{\rm det}
 \xrightarrow{\operatorname{Res}_{\rm det}}\mathcal G_{\rm det}
 \xrightarrow{\operatorname{pr}_{X,{\rm det}}}X_{\chi,{\rm det}}
 \xrightarrow{\mathcal A_{\rm det}}\mathfrak M_{\rm det}
 \xrightarrow{\mathcal P_{\rm det}}\mathcal C_{\rm det}^{\rm HMT}.
\end{equation}
Each arrow has an inverse on its image given by a coordinate
projection. History is not replaced by value, multisection by selector,
complex by determinant or proof by name.

\section{Theorem of Internal Determinantal Realization}

\begin{theorem}[Complete internal determinant realization]
\label{thm:detint-realizacion}
For every \(x\in\mathcal C_{\rm cont}^{\rm det}\), the chain
\eqref{eq:detint-cadena-completa} constructs a macro-object natural in
depth and precision. It contains two cycles with a common root, an explicit filler
that identifies them, a compatible group-like unit, and a
\(3\)-adic terminal whose order, Smith form, and Bockstein tower produce the same
oriented reader. The thirteen coordinates of \eqref{eq:detint-d-trece}
remain recoverable in the output.
\end{theorem}

\begin{proof}
Equations \eqref{eq:detint-p0-diferencial}--
\eqref{eq:detint-d2} prove that \(P_n^0(c)\) is a complex. The two
augmentations are maps of complexes, so the pullback is typed.
Alexander--Whitney is natural and makes each vertex group-like; together with
\(r\), this gives the common unit. The direct calculation
\eqref{eq:detint-zpm-ciclos}--\eqref{eq:detint-filler-identidad} proves that
the publications are cycles, have the same root and are joined by
\(h_*\). The formulas for \(\Psi_\alpha\) prove functoriality, and
\(K_\alpha\) proves homotopic naturality while preserving the memory
increment. The definition of \(\mathcal C_{\rm cont}^{\rm det}\) ensures that
terminal reduction produces a square block without postulating it outside the
domain. Smith over \(\mathbf Z_3\) gives \(d=\sum_i a_i\); the calculation of each
block \(3^{a_i}u_i\) gives
\(\Theta_{\rm Bock}=\overline{\lambda^{\rm or}}\). Truncation, reduction and
addition of unit blocks preserve these identities. Finally, the
three graph constructions retain the input as the first coordinate,
so the complete composition is non-ablative.
\end{proof}

\section{Provenance and Forgers}

\paragraph{Origin.}
The TPK nerve, Alexander--Whitney, the counit, the common unit, the normal
form \(\partial f=3ce+b\), \(\partial e=g\),
\(\partial b=-3cg\), the common root, the filler, and the tower architecture
have the status of \texttt{PREEXISTING\_AUTHORIAL\_ARCHITECTURE}. The literal computation
of the two cycles, their root, and the identity
\(\partial h_*=z_--z_+\) has the status of a \texttt{RECOVERED\_RESULT}. The
local system \(\Psi_\alpha/K_\alpha\), the balanced surviving subdomain,
the exclusively \(3\)-adic terminal, and the packaging by
graphs have the status of a \texttt{NEW\_FORMALIZATION}; they formalize without altering
the conceptual provenance of the apparatus.

\begin{ablation}[Internal falsifiers of the determinantal branch]
\label{abl:detint-falsadores}
The construction is falsified at the corresponding level if any of
the following facts occurs.
\begin{enumerate}
 \item \(\partial^2=0\) fails; then the local complex does not exist.
 \item Alexander--Whitney does not commute with refinement, the counit fails, or
 the vertex ceases to be group-like.
 \item \(z_\pm\) are not cycles, their roots differ, or
 \(\partial h_*=z_--z_+\) fails.
 \item \(\Psi_\alpha\) is not a map of complexes or
 \(K_{\beta\alpha}=K_\beta+\Psi_\beta K_\alpha\) fails; in that case
 carry or memory has been lost.
 \item The reduced complex is not balanced or is not acyclic over
 \(K_3\); the history lies outside the typed terminal domain.
 \item The Bockstein lengths do not coincide with the \(a_i\), or
 \(\Theta_{\rm Bock}\ne\overline{\lambda^{\rm or}}\).
 \item Refinement modifies \(d\) or \(\lambda^{\rm or}\) outside
 unit blocks and recorded orientation changes.
 \item One of the thirteen coordinates is omitted: \emph{object replaced;
 conclusion not transferable}.
\end{enumerate}
These forgers check the internally constructed object; they do not replace it with a subsequent control.
\end{ablation}
\fi
